\documentclass[11pt, a4paper]{article}

\usepackage{fontspec}
\usepackage{geometry}
\usepackage{fancyhdr}
\usepackage[hidelinks]{hyperref}
\usepackage[normalem]{ulem}
\usepackage{multicol}
\usepackage{enumitem}
\usepackage{amsmath}

\geometry{
  top=2cm,
  bottom=2cm,
  left=2cm,
  right=2cm,
  marginparsep=4pt,
  marginparwidth=1cm
}

\renewcommand{\headrulewidth}{0pt}
\pagestyle{fancyplain}
\fancyhf{}
\lfoot{}
\rfoot{}

\setlength{\parindent}{0pt}
\setlength{\parskip}{0pt}

\usepackage{xunicode}
\defaultfontfeatures{Mapping=tex-text}

\setlist[itemize]{leftmargin=1.4em, itemsep=0pt, topsep=1pt, parsep=0pt}

\newcommand{\code}[1]{\texttt{#1}}

\begin{document}

\footnotesize
\sloppy

\begin{center}
\normalsize\textbf{DSST389: Exam II --- Reference Sheet}
\end{center}

\smallskip

\textit{This sheet is attached to your exam. You do not need to memorize any of it.}

\smallskip

\textbf{HTTP requests.}
A request needs an endpoint, a dictionary of parameters, and a \code{User-Agent} header.
The response carries a status code and, usually, JSON.

\begin{itemize}
  \item \code{session.get(url, params=params, headers=headers)} --- send the request.
  \item \code{response.status\_code} \quad \code{response.json()} \quad
        \code{response.raise\_for\_status()}
  \item \textbf{200} success. \quad \textbf{404} not found: a stable answer, worth caching.
        \quad \textbf{429} rate limited: wait and retry. \quad \textbf{5xx} server error:
        retry. Only 200 and 404 go in the cache.
\end{itemize}

\smallskip

\textbf{Pagination.}
A response holds at most a fixed number of results. If more remain, it carries a
\code{continue} block, whose values are merged into the parameters of the next request.
The loop stops when the block is \textit{absent}.

\begin{verbatim}
while True:
    data = polite_get(base_url, params).json()
    ...                                    # collect rows from this response
    if "continue" not in data:
        break
    params.update(data["continue"])
\end{verbatim}

\smallskip

\textbf{Titles, redirects, and identifiers.}
Passing \code{redirects=1} makes the server report every substitution it applied.
\code{normalized} fixes mechanical issues (underscores, the first letter's case);
\code{redirects} follows a real pointer page. They \textit{chain}, so a title is walked
through the map until it stops changing. A title matching no article appears in
\code{pages} under a negative key with a \code{missing} flag and no \code{pageid}.

\begin{itemize}
  \item \code{page\_id} --- the stable key. Survives a page being renamed. Primary key.
  \item \code{doc\_id} --- the readable key. The article title. Unique within the corpus,
        but not stable across renames.
  \item \code{rev\_id} --- names one immutable version of one page, forever.
  \item \textbf{Raw layer}: the requests cache, the exact bytes the server sent.
        \textbf{Processed layer}: the tidy Parquet tables. The second can always be rebuilt
        from the first, without touching the network.
\end{itemize}

\smallskip

\textbf{Dates and datetimes.}

\begin{itemize}
  \item \code{pl.date(year, month, day)} --- build a \code{Date} from three columns.
        \code{pl.datetime(y, m, d, h, m, s)} builds a \code{Datetime}.
  \item \code{c.col.str.to\_datetime()} --- parse a string. Needs no format string if the
        text is ISO-8601 (\code{YYYY-MM-DD HH:MM:SS}); otherwise pass one, e.g.
        \code{"\%Y-\%m-\%dT\%H:\%M:\%SZ"}.
  \item Components: \code{.dt.year()}, \code{.dt.month()}, \code{.dt.day()},
        \code{.dt.hour()}, \code{.dt.ordinal\_day()}, \code{.dt.quarter()},
        \code{.dt.date()}, and \code{.dt.weekday()} (Monday = 1, Sunday = 7).
  \item \code{.dt.truncate("1h")} --- round \textit{down} to a unit, turning a timestamp into
        a bin you can group by. \code{.dt.round()} rounds to the nearest instead.
        Durations: \code{"1h"}, \code{"1d"}, \code{"30d"}, \code{"1mo"}.
  \item \code{.scale\_x\_date(date\_breaks="1 month", date\_labels="\%Y-\%m")} and
        \code{.scale\_x\_datetime(...)} label a time axis.
\end{itemize}

\smallskip

\textbf{Time zones.}
A datetime is either \textit{naive} (no zone) or \textit{aware} (tagged with one). The two
methods below look alike and do opposite things.

\begin{itemize}
  \item \code{.dt.replace\_time\_zone("UTC")} --- \textbf{attaches} a zone to a naive
        datetime. The clock reading does not change. This is a \textit{claim} about what the
        data always meant. Calling it on an already-aware column moves the instant.
  \item \code{.dt.convert\_time\_zone("America/New\_York")} --- \textbf{re-expresses} an aware
        datetime in another zone. The clock reading \textit{does} change. This is a
        \textit{calculation}. Calling it on a naive column silently assumes UTC.
  \item Wikipedia reports every revision timestamp in UTC. Attach, then convert.
\end{itemize}

\smallskip

\textbf{Durations.}
Subtracting two datetimes gives a \code{Duration}. Durations average exactly.

\begin{itemize}
  \item \code{.dt.total\_hours()} \quad \code{.dt.total\_days()} \quad
        \code{.dt.total\_minutes()} \quad \code{.dt.total\_seconds()}
  \item These \textbf{truncate toward zero}, and they do it to each value \textit{before} any
        aggregation. A gap of 5h 35m becomes 5 hours, or 0 days. Choose a unit finer than the
        intervals you are measuring.
\end{itemize}

\textbf{Window functions.}
An ordered function needs a \code{.sort()} in front of it to say what ``previous'' means, and
an \code{.over()} to stop it running across the boundary between one group and the next.

\begin{itemize}
  \item \code{c.size.shift(1).over(c.doc\_id)} --- the value one row back, within each group.
        \code{shift(-1)} looks forward; \code{shift(2)} reaches back two.
  \item \code{c.size.diff().over(c.doc\_id)} --- shorthand for \code{c.size - c.size.shift(1)}.
        The first row of each group is \code{null}.
  \item \code{c.size.rolling\_mean(window\_size=3)} --- a window of \textbf{three rows}. Only
        correct if rows are evenly spaced.
  \item \code{c.size.rolling\_mean\_by(by=c.timestamp, window\_size="30d")} --- a window of
        \textbf{thirty days}, however many rows that is. Use this whenever observations are
        unevenly spaced in time. Also \code{rolling\_sum\_by}, \code{rolling\_max\_by},
        \code{rolling\_std\_by}.
  \item A comparison against \code{null} is \code{null}, not \code{False}, and
        \code{.filter()} keeps only rows that are \code{True}.
\end{itemize}

\smallskip

\textbf{Joins on time.}

\begin{verbatim}
df.join_asof(other, left_on="timestamp", right_on="moment", by="doc_id",
             strategy="backward", tolerance=None)
\end{verbatim}

\begin{itemize}
  \item Matches each left row to the \textit{nearest} right row rather than an equal one. Both
        tables must be sorted on the join key. \code{by=} forces an exact match on a group
        first, playing the same role that \code{.over()} plays for a window function.
  \item \code{strategy}: \code{"backward"} (the default) takes the most recent record at or
        before; \code{"forward"} takes the next at or after; \code{"nearest"} takes whichever
        is closer, and so can reach into the future. An asof join does not return \code{null}
        for want of a \textit{nearby} match, however far it has to reach; \code{tolerance}
        caps the distance.
  \item \code{df.join\_where(other, expr, ...)} --- keep every pair of rows satisfying a
        filter expression. Columns of the right table carry a \code{\_right} suffix.
\end{itemize}

\smallskip

\textbf{The annotation table.}
\code{DSText.process(docs, nlp)} returns one row per token. The primary key is
\code{doc\_id}, \code{sid}, \code{tid} together.

\begin{itemize}
  \item \code{sid} sentence number; \code{tid} token number \textbf{within the sentence};
        \code{token} the text; \code{token\_with\_ws} the text plus trailing whitespace;
        \code{lemma} the normalized form (lowercased, singular, infinitive).
  \item \code{upos} the part of speech: \code{NOUN}, \code{VERB}, \code{ADJ}, \code{ADV},
        \code{PROPN}, \code{DET}, \code{ADP}, \code{NUM}, \code{PUNCT}. \code{tag} is the
        fine-grained, English-specific version.
  \item \code{is\_alpha}, \code{is\_stop}, \code{is\_punct} --- Boolean flags. A number such as
        \code{1847} is \textit{not} alphabetic.
  \item \code{dep} the grammatical relation (\code{nsubj}, \code{dobj}, \code{amod},
        \code{det}, \code{prep}, \code{pobj}, \code{compound}, \code{ROOT});
        \code{head\_idx} the \code{tid} of the token it depends on, in the same
        sentence-level coordinate system as \code{tid}. The root points at itself.
  \item \code{ent\_type} the entity type (\code{PERSON}, \code{ORG}, \code{GPE},
        \code{DATE}, \code{WORK\_OF\_ART}), empty if none; \code{ent\_iob} is \code{B} at the
        first token of an entity, \code{I} inside one, \code{O} outside any.
  \item \code{DSText.kwic(anno, keyword, max\_num=10, window=5)} --- print a keyword-in-context
        display. Prints; does not return a table.
\end{itemize}

\smallskip

\textbf{N-grams and collocations.}
Build a bigram by shifting the lemma column up one within each sentence, then pairing.

\begin{verbatim}
next_token = c.lemma.shift(-1).over([c.doc_id, c.sid])
bigram     = c.lemma + " " + c.next_token
\end{verbatim}

Pointwise mutual information compares a bigram's frequency against what independence would
predict. $P(w)$ is a word's count over the total number of words; $P(w_1, w_2)$ is the
bigram's count over the total number of bigrams. Filter to bigrams occurring at least five
times before ranking, or the top of the list fills with pairs of rare words that happened to
co-occur once.

\begin{equation*}
\text{PMI}(w_1, w_2) \;=\; \log_2 \frac{P(w_1, w_2)}{P(w_1) \cdot P(w_2)}
\end{equation*}

\smallskip

\textbf{TF-IDF.}
\code{DSText.compute\_tfidf(anno, min\_df=0.0, max\_df=1.0)} returns one row per non-zero
(document, lemma) pair. It keeps only alphabetic tokens.

\begin{itemize}
  \item \code{tf\_norm} is the term's share of the tokens \textit{in its document}, not a raw
        count. \code{df} is the number of documents containing the term, and $N$ is the total
        number of documents.
  \item \code{min\_df} and \code{max\_df} are \textbf{proportions} in $[0, 1]$, compared with
        \code{<=} and \code{>=}. \code{max\_df=0.5} drops terms in \textit{more than} half the
        documents; a term in exactly half survives. Because the smoothed idf below never
        reaches zero, \code{max\_df} is what removes ubiquitous vocabulary.
\end{itemize}

\begin{equation*}
\text{tfidf} \;=\; \text{tf} \times \left( \log\!\left( \frac{N + 1}{\text{df} + 1} \right) + 1 \right)
\end{equation*}

The log is natural. The trailing $+1$ puts a floor under the idf: a term appearing in
\textit{every} document gets $\log(1) + 1 = 1$, not $0$, and so keeps a non-zero score.

\smallskip

\textbf{Documents as vectors.}
Each document is a point whose coordinates are its TF-IDF scores, one axis per term; most are
zero, a \textit{sparse} embedding. Euclidean distance is poor here, because a long document
sits far from the origin and length swamps vocabulary. Measure the \textbf{angle} instead,
which ignores length. Scores are never negative, so the angle runs from $0$ (identical
proportions) to $\pi/2 \approx 1.571$ (no shared vocabulary).

\begin{equation*}
d(\mathbf{x}_a, \mathbf{x}_b) \;=\; \arccos\!\left( \frac{\mathbf{x}_a \cdot \mathbf{x}_b}{\lVert \mathbf{x}_a \rVert \, \lVert \mathbf{x}_b \rVert} \right)
\end{equation*}

\begin{itemize}
  \item \code{DSText.compute\_distances(anno, min\_df=0.0, max\_df=1.0)} --- the angle distance
        for every pair. Each unordered pair appears \textit{once}, with \code{doc\_id}
        alphabetically before \code{doc\_id2}.
  \item \code{anno.pipe(umap\_text, doc\_id=c.doc\_id, term\_id=c.lemma, n\_components=2)}
        followed by \code{.predict(full=True)} --- project the corpus down to two dimensions,
        returned as the columns \code{dr0} and \code{dr1}, one row per document. Takes
        \code{min\_df} and \code{max\_df} as well.
  \item Two corpora in different languages have different vocabularies, so their vectors live
        in different spaces and cannot be compared directly. Their \textit{structure} can be.
\end{itemize}

\end{document}
